
\begin{abstract}
대화형 AI 에이전트를 위한 기존 벤치마크는 \textit{단일 제어(single-control)} 환경을 시뮬레이션한다. 여기서는 AI 에이전트만 도구를 사용해 세계와 상호작용할 수 있고, 사용자는 수동적인 정보 제공자로 남는다. 이는 기술 지원처럼 사용자가 (공유된) 세계의 상태를 수정하는 데 능동적으로 참여해야 하는 실제 시나리오와 다르다. 이 간극을 해결하기 위해 우리는 \bench를 소개하며, 네 가지 핵심 기여를 제시한다: 
    1) \decpomdp로 모델링한 새로운 \text{Telecom 이중 제어 도메인}으로, 에이전트와 사용자 모두 도구를 사용해 공유되고 동적인 환경에서 행동하며 에이전트의 조정과 의사소통을 함께 시험한다,
    2) 원자적 구성 요소로부터 다양하고 검증 가능한 과제를 프로그램적으로 생성하여 도메인 포괄성과 제어된 복잡성을 보장하는 \text{조합적 과제 생성기},
    3) 도구와 관찰 가능한 상태로 행동이 제약되어 시뮬레이션 충실도를 높이는, 환경과 긴밀히 결합된 \text{신뢰할 수 있는 사용자 시뮬레이터},
    4) 추론에서 비롯된 오류와 의사소통/조정에서 비롯된 오류를 분리하는 것을 포함한 여러 절제를 통한 \text{에이전트 성능의 세분화된 분석}.
    특히 우리의 실험은 에이전트가 no-user에서 이중 제어로 전환될 때 성능이 크게 하락함을 보여 주며, 사용자를 안내하는 일이 갖는 어려움을 부각한다. 전반적으로 \bench는 효과적으로 추론하면서 동시에 사용자의 행동을 안내해야 하는 에이전트를 위한 통제된 테스트베드를 제공한다.\footnote{코드와 데이터는 다음에서 이용할 수 있다: \href{https://github.com/sierra-research/tau2-bench}{https://github.com/sierra-research/tau2-bench}.}
\end{abstract}
